\section*{第 \ref{ch:b2:continuous-reactors} 章——连续反应器与工业过程}

\begin{solution}{exo:b2:continuous-reactors:1}
$\tau = V/Q_v = 2000/0.50 = \qty{4.0e3}{s}$，即 \qty{1.1}{h}。
\end{solution}

\begin{solution}{exo:b2:continuous-reactors:2}
$k\tau = 1.2 \times 10^{-3} \times 1800 = 2.16$。搅拌釜：$X = 2.16/3.16 = 0.68$；
平推流：$X = 1 - \mathrm e^{-2.16} = 0.88$。
\end{solution}

\begin{solution}{exo:b2:continuous-reactors:3}
平推流：$\tau = \ln 100/0.020 = \qty{230}{s}$。搅拌釜：$X/(1 - X) = 99 = k\tau$，
$\tau = \qty{4950}{s}$，长了二十多倍：转化率较高时，
搅拌釜以出口浓度操作，代价高昂。
\end{solution}

\begin{solution}{exo:b2:continuous-reactors:4}
每升：$120 \times 0.50 = \qty{60}{kJ}$ 的热量被释放，使
热容为 \qty{4.2}{kJ/K} 的混合物升温 $\Delta T_{\text{ad}} = \qty{14}{K}$。冷却失效
会使混合物明显升温，但不至于危险，除非溶剂
接近其沸点，或第二个反应在该温度范围内开始。
\end{solution}

\begin{solution}{exo:b2:continuous-reactors:5}
搅拌釜：$Da = X/(1 - X)^2 = 0.8/0.04 = 20$，$\tau = 20/0.010 = \qty{2000}{s}$。
平推流：$Da = X/(1 - X) = 4$，$\tau = \qty{400}{s}$。比值为 5，一级反应时则为 $4/\ln 5 =
2.5$。级数越高，搅拌釜
因浓度低而受损越重。
\end{solution}

\begin{solution}{exo:b2:continuous-reactors:6}
$k\tau = 5.0$。三个串联搅拌釜：$X = 1 - (1 + 5/3)^{-3} = 0.947$。单个搅拌釜：$5/6 =
0.833$。平推流：$1 - \mathrm e^{-5} = 0.993$。
\end{solution}

\begin{solution}{exo:b2:continuous-reactors:7}
$X = 0.70$ 意味着 $k\tau = 0.70/0.30 = 2.33$。流量加倍则 $\tau$ 减半：
$k\tau = 1.17$，$X = 0.54$。要使转化率回到 70\,\%，停留时间
须加倍，体积也随之加倍。
\end{solution}

\begin{solution}{exo:b2:continuous-reactors:8}
释放的热量：$80 \times 2.0 \times 1.0 \times 0.90 = \qty{144}{kW}$。出口物流带走的
热量（自 \qty{330}{K} 升至 \qty{350}{K}）：$4.0 \times 1.0 \times 20 =
\qty{80}{kW}$。夹套需移走其差值，即 \qty{64}{kW}。
\end{solution}

\begin{solution}{exo:b2:continuous-reactors:9}
线性尺寸 $\times 10$：体积 $\times 1000$，器壁面积 $\times 100$，
比表面积 $\div 10$。放出的热量随体积增长，经器壁移走的热量
随面积增长：大型反应器在相同温差下，每升
移走的热量少了十倍。
\end{solution}

\begin{solution}{exo:b2:continuous-reactors:10}
平推流：$\tau = \ln(0.05/0.10)/(0.05 - 0.10) = \qty{13.9}{min}$，$c_B/c_{A,0} =
\frac{0.10}{-0.05}(\mathrm e^{-1.386} - \mathrm e^{-0.693}) = 0.50$。搅拌釜：$\tau =
1/\sqrt{0.005} = \qty{14.1}{min}$，$c_B/c_{A,0} = 1.414/(2.414 \times 1.707) =
0.34$。搅拌釜中，新鲜的 A 与已生成的 B 在同一体积内混合：总有一部分
B 停留足够久而转化为 C，同时有一部分 A 刚到达；
平推流中，每一薄层流体经历的历程都相同。
\end{solution}

\begin{solution}{exo:b2:continuous-reactors:11}
提高冷却剂（及进料）温度，使移热线右移。
冷态交点缓慢上升，直至移热线与
S 形曲线的下拐点相切；此后冷态消失，釜内
跃迁至热态分支（点火）。再次降温时，釜内仍
停留在热态分支，直至移热线与上拐点相切，才
回落至冷态（熄火），其温度低于点火温度，形成
磁滞回线。
\end{solution}

\begin{solution}{exo:b2:continuous-reactors:12}
间歇反应器：\qty{2.0}{mol/L} 中仍有 75\,\% 未反应，$1.5 \times 100 = \qty{150}{kJ/L}$，
$\Delta T = 150/4.0 = \qty{37.5}{K}$，超出操作温度，足以
达到沸点或发生分解。半间歇操作时，未反应物至多
\qty{0.10}{mol/L}，$\Delta T = 10/4.0 = \qty{2.5}{K}$。缓慢加料
可限制反应器中积蓄的能量。
\end{solution}

\begin{solution}{pb:b2:continuous-reactors:1}
\textbf{1.} $r = k[\text{酯}][\ce{OH-}]$；等量进料且
化学计量比为 1 : 1，使两者浓度保持相等，$r = kc^2$。
\textbf{2.} $1/c - 1/c_0 = kt$，当 $c = 0.10c_0$ 时：$kc_0t = 9$，$t = 9/(0.11
\times 0.10) = \qty{818}{s}$，约 \qty{14}{min}。
\textbf{3.} $t_{1/2} = 1/(kc_0) = \qty{91}{s}$；对于二级反应，速率随
浓度的平方而下降，因此使其减半所需的时间取决于
起始浓度。
\textbf{4.} $Da = kc_0\tau = 0.011\,\tau$（$\tau$ 以秒为单位）。
\textbf{5.} $Q_vc_0 - Q_vc - kc^2V = 0$，即 $c_0 - c = k\tau c^2$。
\textbf{6.} 代入 $c = c_0(1 - X)$：$c_0X = k\tau c_0^2(1 - X)^2$。
\textbf{7.} $Da = 0.90/0.10^2 = 90$。
\textbf{8.} $\tau = 90/0.011 = \qty{8.2e3}{s}$，即 \qty{2.3}{h}。
\textbf{9.} $V = 0.50 \times 8182 = \qty{4.1e3}{L}$，即 \qty{4.1}{m^3}。
\textbf{10.} 出口浓度为 $0.010\,\unit{mol/L}$：各处的速率
均比入口低一百倍。
\textbf{11.} $Da = X/(1 - X) = 9$，$\tau = \qty{818}{s}$，$V = \qty{409}{L}$。
\textbf{12.} 平推流所需的体积少了十倍。
\textbf{13.} $c_{j-1} - c_j = k\tau_1c_j^2$，其中 $\tau_1$ 为
单个搅拌釜的停留时间。
\textbf{14.} $\tau_1 = 850/(3 \times 0.50) = \qty{567}{s}$，$k\tau_1 =
\qty{62.4}{L/mol}$。由 $62.4c_j^2 + c_j - c_{j-1} = 0$ 解得：$c_1 = 0.0328$，
$c_2 = 0.0163$，$c_3 = \qty{0.0100}{mol/L}$，转化率为 90\,\%。
\textbf{15.} 平推流，或串联搅拌釜（二者体积分别比单个搅拌釜
小十倍和五倍）。单个搅拌釜仍可能占优，因其温度均匀、
便于清洗和控制，且此处体积成本低。
\textbf{16.} $\Delta T_{\text{ad}} = 55\,000 \times 100/(997 \times 4180) =
\qty{1.3}{K}$（转化率 90\,\% 时为 \qty{1.2}{K}）。
\textbf{17.} $55 \times 0.50 \times 0.10 \times 0.90 = \qty{2.5}{kW}$。
\textbf{18.} 否：出口物流带走热量，温升约为
一开尔文，不可能发生热失控。
\textbf{19.} 浓度提高二十倍：绝热温升变为
\qty{26}{K}，放热功率 \qty{50}{kW}，值得加装冷却盘管；速率 $kc_0$
提高了二十倍，故达 90\,\% 转化率所需的釜体积缩小二十倍
（\qty{0.20}{m^3}）。
\textbf{20.} $0.50 \times 0.10 \times 0.90 = \qty{0.045}{mol/s}$，$\times 86\,400
\times 82.03 = \qty{319}{kg}$ 乙酸钠（每天）。
\textbf{21.} 速率常数随温度升高而增大（阿伦尼乌斯），因此
体积随之缩小；代价是进料加热，而对于其他
反应则会导致选择性降低、蒸气增多。
\textbf{22.} 单个搅拌釜须具有 $\boldsymbol{V \approx
\qty{4.1}{m^3}}$。
\end{solution}
